\documentclass{ximera}

\input{../../preamble.tex}

\title{Rectangular}

\begin{document}

\begin{abstract}
coordinates
\end{abstract}
\maketitle




The Cartesian plane graphs dots using rectangular coordinates.

\begin{center}
(Left/Right, Up/Down)
\end{center}

These coordinates give the directional measurements from the origin to the point's location.

We then extend this for functions.


We use the coordinates of plotted points to visually encode function information. A function is a collection of domain-range pairs. We can view these pairs as ordered pairs and interpret the domain number as left/right information and interpret the funciton value as up/down information.  We cna then plot the point representing the domain-function value pair.

Let $f(x) = 3 \, \sin(2x+\tfrac{\pi}{4}) - 1$

\begin{image}
\begin{tikzpicture}
  \begin{axis}[
            domain=-10:10, ymax=5, xmax=10, ymin=-5, xmin=-10,
            axis lines =center, xlabel={$\theta$}, ylabel=$y$, grid = major, grid style={dashed},
            ytick={-4,-3,-2,-1,1,2,3,4},
            xtick={-7.85, -6.28, -4.71, -3.14, -1.57, 0, 1.57, 3.142, 4.71, 6.28, 7.85},
            xticklabels={$-\tfrac{5\pi}{2}$,$-2\pi$,$-\tfrac{3\pi}{2}$,$-\pi$, $-\tfrac{\pi}{2}$, $0$, $\tfrac{\pi}{2}$, $\pi$, $\tfrac{3\pi}{2}$, $2\pi$, $\tfrac{5\pi}{2}$},
            yticklabels={$-3$,$-2$,$-1$,$1$,$2$,$3$,$4$}, 
            ticklabel style={font=\scriptsize},
            every axis y label/.style={at=(current axis.above origin),anchor=south},
            every axis x label/.style={at=(current axis.right of origin),anchor=west},
            axis on top
          ]
          
            \addplot [line width=2, penColor2, smooth,samples=300,domain=(-10:10),<->] {3*sin(deg(2*x+0.785))-1};



  \end{axis}
\end{tikzpicture}
\end{image}


Each point on the graph visually encodes a domain-range pair in the function $f$

\[  (a, f(a)) = (a, 3 \, \sin(2a+\tfrac{\pi}{4}) - 1)   \]




We can view this function as a composition of the core sine function with two linear functions.



\begin{itemize}
\item Let $L_{out}(u) = 3 \, u - 1$, an increasing linear function
\item Let $S(\theta) = \sin(\theta)$, the core sine function
\item Let $L_{in}(v) = 2 \, v + \tfrac{\pi}{4}$, an increasing linear function
\end{itemize}



\[ 3 \, \sin(2x + \tfrac{\pi}{4}) - 1  = (L_{out} \circ S \circ L_{in})(x)\]



When viewed this way, we say that the function $f$ is a \texbf{transformation} of $\sin(\theta)$.


There is a domain transformation given by $L_{in}(v) = 2 \, v + \tfrac{\pi}{4}$.

There is a range transformation given by $L_{out}(u) = 3 \, u - 1$.





\subsection*{Graphical Transformation}



Along with the function, we say that there has been a \textbf{graphical transformation} as well.



There is a horizontal transformation given by $L_{in}(v) = 2 \, v + \tfrac{\pi}{4}$.

There is a vertical transformation given by $L_{out}(u) = 3 \, u - 1$.




The transformations are rectangular, because the measurements on the graph are rectangular.  Domain is measured horizontally. Range (function value) is measured vertically.




\textbf{\textcolor{blue!55!black}{Horizontal Transformations}} 

Our transformations are \textbf{linear}, composed of shifts and stretches.

We can describe the shifts and stretches by following the arithmetic performed on the old variable to obtain the new variable.

Horizontally, $\sin(\theta)$ is transformed into $3 \, \sin(2x+\tfrac{\pi}{4}) - 1$, by replacing $\theta$ with $2x+\tfrac{\pi}{4}$.


\[ \theta =  2x+\frac{\pi}{4}\]

But, this shows the arithmetic to perform on $x$, the new variable.  We need to solve for 'x' in terms of '$\theta$''.



 \[ \frac{1}{2} \left( \theta -  \frac{\pi}{4} \right) =  x\]


First, subtract $\frac{\pi}{4}$, which is a horizontal shift left.

Second, multiply by $\frac{1}{2}$, which is a horzontal compression.






\textbf{\textcolor{blue!55!black}{Vertically Transformations}} 



Vertically, $\sin(\theta)$ is transformed into $3 \, \sin(2x+\tfrac{\pi}{4}) - 1$, by replacing $\sin$ with $3 \, \sin - 1$.

This shows exact what to do with the old function to get the new function, so there is nothing to be done.

The old function is first multiplied by $3$, which is a vertical stretch.

Then $1$ is subtracted, whcih is a vertical shift down.






\begin{image}
\begin{tikzpicture}
  \begin{axis}[
            domain=-10:10, ymax=5, xmax=10, ymin=-5, xmin=-10,
            axis lines =center, xlabel={$\theta$}, ylabel=$y$, grid = major, grid style={dashed},
            ytick={-4,-3,-2,-1,1,2,3,4},
            xtick={-7.85, -6.28, -4.71, -3.14, -1.57, 0, 1.57, 3.142, 4.71, 6.28, 7.85},
            xticklabels={$-\tfrac{5\pi}{2}$,$-2\pi$,$-\tfrac{3\pi}{2}$,$-\pi$, $-\tfrac{\pi}{2}$, $0$, $\tfrac{\pi}{2}$, $\pi$, $\tfrac{3\pi}{2}$, $2\pi$, $\tfrac{5\pi}{2}$},
            yticklabels={$-3$,$-2$,$-1$,$1$,$2$,$3$,$4$}, 
            ticklabel style={font=\scriptsize},
            every axis y label/.style={at=(current axis.above origin),anchor=south},
            every axis x label/.style={at=(current axis.right of origin),anchor=west},
            axis on top
          ]
          
            \addplot [line width=2, penColor2, smooth,samples=300,domain=(-10:10),<->] {3*sin(deg(2*x+0.785))-1};

             \addplot [line width=2, penColor, smooth,samples=300,domain=(-10:10),<->] {sin(deg(2*x))};




  \end{axis}
\end{tikzpicture}
\end{image}
































\begin{center}
\textbf{\textcolor{green!50!black}{ooooo-=-=-=-ooOoo-=-=-=-ooooo}} \\

more examples can be found by following this link\\ \link[More Examples of Trigonometric Graphs]{https://ximera.osu.edu/csccmathematics/precalculus2/precalculus2/trigonometricGraphs/examples/exampleList}

\end{center}





\end{document}

